\documentclass[11pt]{article}
\usepackage[margin=1in]{geometry}
\usepackage{amsmath,amssymb}
\newcommand{\finalanswer}[1]{\par\medskip\noindent\fbox{\parbox{0.94\linewidth}{\textbf{Answer:} #1}}}
\title{STAT 412 Homework 5 --- Question 1}
\author{}\date{}
\begin{document}\maketitle
\section*{Problem}
Find the mean, median, and mode for the sample of sick days:
\[
9,\ 16,\ 13,\ 10,\ 4,\ 68,\ 20,\ 9,\ 14.
\]
Which value appears to be the best measure of the center? Complete the
comparison sentence:
\[
\text{Since the mean is \_\_\_ the median is \_\_\_ and the mode is \_\_\_
most of the data values, the \_\_\_ appears to be the best measure of center.}
\]

\section*{Work}
\[
\bar x=\frac{9+16+13+10+4+68+20+9+14}{9}
=\frac{163}{9}=18.11.
\]
Sorted data:
\[
4,\ 9,\ 9,\ 10,\ 13,\ 14,\ 16,\ 20,\ 68.
\]
The middle value is $13$, so the median is $13$. The value occurring most often
is $9$, so the mode is $9$ (not $68$; $68$ occurs only once).

Since $68$ is a large high outlier, it pulls the mean upward. Thus the mean
$18.11$ is greater than the median $13$. The mode $9$ is less than most of the
data values, so it is also not the best summary of the center. The median is
the best measure here because it is resistant to the high outlier.

\finalanswer{Mean $=18.11$; median $=13$; mode $=9$. Dropdowns:
\emph{greater than}, \emph{less than}, \emph{median}.}
\end{document}
